% Folder 136, batch 18 — archivists' pages 341–360 (pages 1–20 of batch-18.pdf).
% Pass: Opus 5.5 (claude-opus-5-5), 2026-10-03 — first pass, unchecked against the pages by a human.
%
% Page numbering: \page{} follows the batch numbering (batch page 1 = 341). The
% pencilled archive numbers visible on the leaves read two higher (343 on batch
% page 1, … 362 on batch page 20); the facsimile is turned by the batch numbering.
% Skipped: 343, 346, 353, 355 (blank leaves); 348, 350, 352 (typescript by another
% hand, not transcribed — RIGHTS.md). Pages 347, 349 and 351 are in very faint
% pencil; only fragments are read.
\documentclass[11pt,a4paper]{article}
\input{../preamble/grothendieck}

\folder{136}
\batch{18}
\pages{341}{360}
\foldertitle{Complexe de De Rham à puissance divisée [conférence de 1976 à l’IHÉS] : notes manuscrites (s.d.).}
\dating{[à partir de 1975-1976]}
\watermark{Édition de démonstration}

\begin{document}

\section{Le foncteur $K$ et les complexes $\Phi_*$, $\Psi_*$}

\page{341}
\note{la numérotation au crayon des archivistes, sur les feuillets de ce lot, est supérieure de deux à celle qu’emploie cette édition (ce feuillet porte « 343 ») ; les renvois de Grothendieck, s’il y en a, sont à lire avec ce décalage. Le calcul de la page reprend des notations ($F$, $F'$, $K$, $M_*$, $\Phi_*$, $\Psi_*$) fixées avant ce lot.}

\[
K \longrightarrow F \rightrightarrows F' ,
\qquad F \nearrow G \qquad (G = F_{\Delta^n_*})
\]
\[
K(M_*) \longrightarrow F(M_*) \rightrightarrows F'(M_*)
\]
\struck{$C(M_*)$}

\[
\Delta^0_* \rightrightarrows \Delta^1_* \xrightarrow{\;S\;} \cdots
\]
\[
K(M_*) \longrightarrow M_{[1]} \overset{\alpha^*_{M_*}}{\underset{\beta^*_{M_*}}{\rightrightarrows}} M_{[0]} ,
\qquad K(M_*) = F_S(K)
\]
\uncertain{fonctoriel ?} — la flèche $u_M \colon K(M_*) \to M_{[n]}$.
\note{sous $K(M_*)$ on lit un « $F''_S(K)$ » rattaché par un signe d’égalité vertical ; la lecture de l’accent est douteuse.}

\[
F_S \longrightarrow F_{\Delta^1_*} \rightrightarrows F_{\Delta^0_*} ,
\qquad
k \rightrightarrows \Lambda^2 k^2 \simeq k \longrightarrow 0
\]
\[
\uparrow \wedge T_2 , \qquad \uparrow i_0 , \qquad
k \longrightarrow k^2 \rightrightarrows k \quad (\alpha \mapsto 0)
\]
\[
\alpha T_n \wedge \beta \in \Lambda^2 \Phi_{[n]}
\quad \uparrow \quad
\alpha T_{[n]} \in k^{\Delta^n} \simeq \Phi_{[n]} ,
\qquad \alpha \in k , \ \beta \in \Psi_{[n]}
\]
\struck{$\mathrm{Hom}(\alpha^*_M , \beta^*_M)$}

\[
K(M_*) \otimes K(N_*) \longrightarrow M_{[1]} \otimes N_{[1]} \rightrightarrows M_{[0]} \otimes N_{[0]}
\]
\[
u_{M_*} \otimes u_{N_*} \colon K(M_*) \otimes K(N_*) \longrightarrow M_{[n]} \otimes N_{[n]}
\]
\[
K(M_* \otimes N_*) \longleftarrow [M \otimes N]_1 \rightrightarrows [M \otimes N]_0 ,
\qquad
u_{M_* \otimes N_*} \colon K(M_* \otimes N_*) \longrightarrow (M_* \otimes N_*)_{[n]}
\]
\note{les flèches verticales qui relient la ligne des $M_{[i]} \otimes N_{[i]}$ à celle des $[M \otimes N]_i$ portent un signe lu « 2 », peut-être un $\wr$ ; après $(M_* \otimes N_*)_{[n]}$ un « $= M_{[n]} \otimes N$ » est biffé.}

À droite :
\[
\begin{aligned}
K(M_*) &\simeq Z^0(M_\bullet \otimes K^\bullet) \simeq \mathrm{Hom}^\bullet(M_\bullet^\vee , K^\bullet) \\
&\simeq \mathrm{Hom}_*(M_*^\vee , K^*) \simeq \mathrm{Hom}_*(K_* , M_*)
\end{aligned}
\]
\[
K^*_* \simeq \Psi_* \oplus \Lambda^2 \Psi_* =
\]

Prenons $\Phi_* = C^0_*$.
\[
K(\Phi_*) \simeq k \longrightarrow F_{\Delta^n_*}(\Phi_*) \simeq \Phi_{[n]} = k^{\Delta^n_*}
\]
\note{l’indice de l’exposant $\Delta^n_*$ est peu lisible.}
\struck{$\mathrm{Hom}($} \quad \struck{$K(M_*)$} \quad \struck{$\mathrm{Hom}($}
\[
\mathrm{Hom}_*(M_* , K_*)
\]
\[
K_M(K) = \mathrm{Hom}(M , K)
\]
\[
\mathrm{Hom}(M , K) \otimes \mathrm{Hom}(M , L) \longrightarrow \mathrm{Hom}(M , K \otimes L) ,
\qquad M \longrightarrow M \otimes M
\]

\page{342}
\note{page de calculs dispersés ; l’ordre de lecture adopté va de l’encadré du haut vers le bas, puis la colonne de droite.}

\uncertain{Encadré} : $A \xrightarrow{\;u_i\;} B$, \ill{} \uncertain{injectif} \ill{}.
\[
\textstyle\bigoplus k[T] \xrightarrow{\;u_i\;} B , \qquad T \mapsto b_i , \quad T^n \mapsto b_i^n ,
\qquad \sum \lambda_i b_i^n = 0
\]
\[
\textstyle\sum \lambda_i u_i = 0 \;\Rightarrow\; \sum \lambda_i u_i(x)^n = 0 \quad \forall n
\]
$\mathrm{Ker}\, u_i = \mathfrak{p}_i$ idéaux premiers \ill{} ; $t \in \bigcup \mathfrak{p}_i =$ \uncertain{$\bigcap \mathfrak{p}_i$}.
\struck{$u_i = u_j \to u_i(x) - u_j(y)$}

\[
\mathrm{Hom}(\Delta^m_* , \Delta^n_*) \xrightarrow{\;\sim\;} \mathrm{Hom}_{\mathrm{mult}}(F(\Delta^n_*) , F(\Delta^m_*))
\]
\[
X_* \longmapsto \widetilde{X}_* = \bigl( n \mapsto \mathrm{Hom}(F_{X_*} , F_{\Delta^n_*}) \bigr)
\qquad \text{\struck{$n \mapsto X_n$}}
\]
\[
\begin{aligned}
\widetilde{X}_n &\simeq \mathrm{Hom}\bigl( \varprojlim_{\Delta^m_*/X_*} F_{\Delta^m_*} , F_{\Delta^n} \bigr) \\
&\uparrow \text{ iso ?} \\
&\varinjlim_{\Delta^m_*/X_*} \mathrm{Hom}(F_{\Delta^m_*} , F_{\Delta^n})
\simeq \varinjlim \mathrm{Hom}(\Delta^n_* , \Delta^m_*) \simeq \varinjlim \Delta^m_{[n]} \simeq X_{[n]}
\end{aligned}
\]
\note{le « iso ? » est entouré sur la page ; la première ligne n’a pas de membre de gauche écrit, $\widetilde{X}_n$ est suppléé d’après la ligne précédente.}

Soient \struck{donc} $\prod F_i \to F_{\Delta^n}$, avec $m_i \colon F_i \to$ \ill{} et les $u_i$ en diagonale ;
\[
\textstyle\prod G_i \longrightarrow C^\bullet_{[n]}
\qquad \text{\struck{$\sum \lambda_i u_i$}}
\]
$\forall M_*$, \struck{\ill{} $(M_*) \to F(M_*)$}
\[
\mathrm{Sym}\, G(M_*) \xrightarrow{\;u_i\;} \mathrm{Sym}\, F(M_*)
\]

\page{344}
\struck{$F(X_*) =$}

$C^\bullet(X_*)$ muni d’une structure multiplicative et d’une \uncertain{dérivation} \ill{}.
\[
F_*(X_*) = \bigl[ n \mapsto \mathrm{Hom}(C^\bullet(X_*) , C^\bullet(\Delta^n_*)) \bigr]
\]
\[
F_*(\Delta^m_*) = P^m_* \qquad \text{torseur sous } \Psi^\vee_m \otimes \Phi_* = \Phi^m_*
\]
\note{après « torseur sous » une parenthèse, \uncertain{dont le groupe des sections est $k^{m+1}$}, se lit mal.}
\[
0 \longrightarrow \Psi^\vee_m \otimes \Phi_* \longrightarrow \Phi^\vee_m \otimes \Phi_* \longrightarrow \Phi_* \longrightarrow 0 ,
\qquad \mathrm{Hom}(\Phi_m , \Phi_*) , \quad T_m \mapsto T_*
\]
\[
\Delta^m_* \longrightarrow P^m_* , \quad y \in P^m_n ;
\qquad \Delta^m_n = \mathrm{Hom}_{\mathrm{ens}}(\Delta^n , \Delta^m) \longrightarrow P^m_n
\]
$\longrightarrow \mathrm{Hom}(P^m_* , P^n_*)$ \qquad $\longrightarrow \mathrm{Hom}(P^n$ \ill{}

\[
\begin{aligned}
F_*(X_*) &= \mathrm{Hom}\bigl( \varprojlim_{\Delta^m_* \in \mathrm{Ob}\, \Delta/X_*} \Box(C^\bullet(\Delta^m_*)) , C^\bullet(\Delta^n_*) \bigr) \\
&\uparrow \text{ est-ce une équiv.\ d’homotopie ?} \\
&\varinjlim_{\Delta^m \in \mathrm{Ob}\, \Delta/X_*} \mathrm{Hom}\bigl( \Box(C^\bullet(\Delta^m_*)) , C^\bullet(\Delta^n_*) \bigr) = \varinjlim P^m_*
\end{aligned}
\]
\note{le signe noté $\Box$ dans ces deux lignes est \ill{} ; c’est un foncteur appliqué à $C^\bullet(\Delta^m_*)$.}
\[
\uparrow \ \varinjlim \Delta^m_* \quad (\Delta^m \in \mathrm{Ob}\, \Delta)
\]
\[
\Phi_m \longrightarrow \Phi_n , \quad T_m \mapsto T_n , \qquad \alpha \in k , \ \beta \in T_m
\]
\[
T_0 , \quad T_1 \circ T_1 ; \qquad y , \quad T_1 \circ y , \quad y \circ T_1 , \quad y \circ y
\]

\[
X_* = \mathrm{Coker}\bigl( \Delta^0_* \overset{\partial_0}{\underset{\partial_1}{\rightrightarrows}} \Delta^1_* \bigr) ,
\qquad
C^\bullet(X_*) = \mathrm{Ker}\bigl( C^\bullet(\Delta^1_*) \overset{C(\partial_0)}{\underset{C(\partial_1)}{\rightrightarrows}} C^\bullet(\Delta^0_*) \bigr)
\]
\[
k \xrightarrow{\;\alpha \mapsto 0\;} \Phi_1 = k^2 \overset{\partial_0 = \mathrm{pr}_0}{\underset{\partial_1 = \mathrm{pr}_1}{\rightrightarrows}} \Phi_0 = k ,
\qquad
k \xrightarrow{\;\mathrm{id}\;} \Lambda^2 \Phi_1 = k
\]
\[
1 = T_1 \longmapsto \alpha T_m \in \Phi_m ;
\qquad
1 \longmapsto T_m \wedge \beta \in \Lambda^2 \Phi_m \quad (\beta \in \Phi_m)
\]
\[
F_*(X_*) = \mathrm{Hom}\bigl( [\, 0 \to k \xrightarrow{0} k \to 0 \to \cdots ] , C^\bullet(\Delta^n_*) \bigr)
\]
(structure mult.\ induite par $C^\bullet(\Delta^1_*)$)
\[
= \bigl\{ n \mapsto \mathrm{Hom}\bigl( [\cdots] , [\, \Phi_n \to \Lambda^2 \Phi_n \to \cdots ] \bigr) \bigr\}
\]

\page{345}
\note{feuillet en largeur, couvert de calculs serrés ; les relations de commutation sont transcrites dans l’ordre où elles se lisent, colonne de gauche puis colonne de droite. Les astérisques sont de lui et semblent marquer les relations vérifiées.}

Opérateurs $d$, $\lambda = \bar{T}\cdot$ \uncertain{(lecture de la barre)}, $\delta$, $\mu_i = T^{(i)}\cdot$, $\Theta(x) = \deg(x)\, x$ ; éléments $1, T, \bar{T}$.
\struck{$\Lambda\Phi \otimes \Lambda\Psi$}
\[
\begin{aligned}
& d\lambda + \lambda d = 0 \ (*) \\
& d\delta + \delta d = \Theta \\
& \delta\lambda + \lambda\delta = \mu \\
& d\mu_i - \mu_i d = \mu_{i-1} \lambda \\
& \lambda\mu_i - \mu_i\lambda = 0 \ (*) \\
& \delta\mu_i - \mu_i\delta = 0 \ (*)
\end{aligned}
\qquad
\begin{aligned}
& d(T^{(i)} x) = d(T^{(i)})\, x + T^{(i)}\, dx \\
& \delta(T x) = \delta(T)\, x + T\, \delta x \\
& d\Theta - \Theta d = 0 \ (*) , \quad \delta\Theta - \Theta\delta = 0 \ (*) \\
& \Theta\lambda - \lambda\Theta = \lambda \\
& \Theta\mu_i - \mu_i\Theta = i\mu_i
\end{aligned}
\]
\note{la deuxième relation est suivie d’une précision « en degré » \ill{} ; des fragments biffés accompagnent la relation $d\mu_i - \mu_i d = \mu_{i-1}\lambda$.}

Entouré :
\[
\Theta\mu_f - \mu_f\Theta = \mu_{\Theta f} , \qquad d\mu_f - \mu_f d = \mu_{f'} \lambda
\]
\[
d^2 = \delta^2 = \lambda^2 = 0 , \qquad \mu_f \mu_g = \mu_{fg} \ ?
\]
\[
\Theta(\lambda x) = (\deg x + 1)\, \lambda x , \qquad \lambda(\Theta x) = \deg x \; \lambda x ,
\qquad \Theta\lambda - \lambda\Theta = \lambda
\]
\note{plusieurs petits schémas reliant $d, \delta, \lambda, \mu_i, \Theta$ et un décompte $\frac{n(n-1)}{2}$ \ill{} ne sont pas transcrits.}

Tout élément de l’algèbre engendrée s’écrit
\[
\sum \mu_{f_{\alpha\beta\sigma n}} \, \lambda^\alpha d^\beta \delta^\sigma \Theta^n ,
\qquad \alpha, \beta, \sigma \in \mathbb{F}_2 , \ n \in \mathbb{N}
\]
\[
\mu_f \lambda^\alpha d^\beta \delta^\sigma \Theta^n \cdot \mu_g \lambda^{\alpha'} d^{\beta'} \delta^{\sigma'} \Theta^{n'} ,
\qquad d^\beta \mu_g = \mu_g d^\beta + \beta\, \mu_{g'} \lambda
\]
Encadré, la liste des relations :
\[
\begin{aligned}
& d\delta + \delta d = \Theta , \qquad \delta\lambda + \lambda\delta = \mu_T , \\
& d\mu_f - \mu_f d = \mu_{f'}\lambda , \qquad \Theta\lambda - \lambda\Theta = \lambda , \\
& \Theta\mu_f - \mu_f\Theta = \mu_{\Theta f}
\end{aligned}
\]
\struck{$\delta(T^{(i)}) = 0$} \quad \uncertain{$\delta(\bar{T}) = 1$}

Colonne de droite : \uncertain{Cas $p$ quelconque, $q = 1$}.
\[
0 \longrightarrow \Gamma^p \Phi_* \longrightarrow \Gamma^{p-1}\Phi_* \otimes \Phi_* \longrightarrow \Gamma^{p-1}\Phi_* \otimes \Psi_* \longrightarrow 0
\]
\[
\Phi_* \times \Psi_* \longrightarrow \ \cdot\ , \qquad (\varphi , \psi) \longmapsto s(\varphi \otimes \psi)\, \varphi^{(p-1)}
\]
\[
\Gamma^{*+1}\Phi_* \otimes \Lambda^* \Psi_* \ \downarrow \ \Gamma^{*+1}\Phi_* \otimes \Lambda^{*}\Psi_*
\]
\note{le dernier facteur de la suite exacte ($\Psi_*$) et la ligne $\Gamma^{*+1}\Phi_* \otimes \Lambda^* \Psi_*$ sont de lecture douteuse.}
\[
\varphi_1^{(p_1)} \cdots \varphi_r^{(p_r)} \longmapsto \cdots , \qquad \sum p_i = p
\]
\[
f\Bigl( \sum^r \lambda_i \varphi_i \Bigr) = s\Bigl( \sum \lambda_i \varphi_i \otimes \psi \Bigr) \Bigl( \sum \lambda_i \varphi_i \Bigr)^{(p-1)}
\]
\struck{$\sum \lambda_i \varphi_i$}

\section{Algèbres différentielles sur un topos}

\page{347}
\note{crayon très pâle ; seuls les mots et formules lisibles sont donnés, le reste est marqué \ill{}.}

$k$-Algèbres différentielles $C^*$ sur \uncertain{Topos} $\mathcal{X}$.

Anneau de base $k$ (\ill{}).
\uncertain{résolution} $C^*$ de $k$ :
\[
\cdots \longrightarrow C^0 \longrightarrow C^1 \longrightarrow C^2 \longrightarrow \cdots ,
\qquad k \longrightarrow C^0
\]
\marginal{\ill{} \uncertain{augmentation}}
(\ill{} \uncertain{correspond à une section $T$ de $C^0$}).

\uncertain{Accouplements} (\uncertain{notés} $x, y \mapsto x \circ y$)
\[
\begin{cases}
Z^i \times C^j \longrightarrow C^{i+j} & (a_{ij}) \\
C^i \times Z^j \longrightarrow C^{i+j} & (b_{ij})
\end{cases}
\qquad i, j \ge 0
\]
satisfaisant aux \uncertain{conditions} :

1°) $a_{ij}$ et $b_{ij}$ \uncertain{coïncident} sur $Z^i \times Z^j$ \ill{} \uncertain{y prennent leurs valeurs dans} $Z^{i+j}$. (\ill{} la notation $x \circ y$ \uncertain{non ambiguë}, \ill{} $dx \circ dy$ \ill{}.)

2°) Associativité : $(x \circ y) \circ z = x \circ (y \circ z)$ \ill{} \uncertain{des trois sections locales} $x, y, z$ (\ill{} $C^i, C^j, C^k$) \ill{} des cocycles ; $1 \circ x = x \circ 1 = x$.

3°)
\[
\begin{cases}
x \circ y = y \circ x \, (-1)^{\deg y \deg x} \\
x \circ x = 0 \quad \text{si } x \text{ section locale de } Z^i
\end{cases}
\]
\ill{} (condition sur $i$) \ill{} \uncertain{bilinéaires} \ill{}.

4°) \uncertain{Pour tout} $i, j \ge 1$, \uncertain{existe} \ill{}
\struck{$C^{i-1} \times Z^j \xrightarrow{\;a_{ij}\;} C^{i+j-1}$}
\[
\begin{array}{ccc}
Z^{i} \times Z^j & \longrightarrow & Z^{i+j} \\
\uparrow{\scriptstyle d \times \mathrm{id}} & & \uparrow{\scriptstyle d} \\
C^{i-1} \times Z^j & \xrightarrow{\;b_{ij}\;} & C^{i+j-1}
\end{array}
\]
\ill{}
\note{la lecture du carré est incertaine : les exposants et l’étiquette de la flèche du bas sont pâles ; une ligne biffée suit, où se lit $d(x^{(i)}) \times Z^j = d($\ill{}$)$.}
(N.B. \ill{} $\mathrm{Hom}_k($\ill{}$\otimes Z^j , Z^{i+j-1})$ \ill{})
\note{le bas de la page, barré de longs traits obliques, n’est pas lisible.}

\page{349}
\note{crayon très pâle, comme p.~347 ; lecture fragmentaire.}

Sous ces conditions, pour tout $U \in \mathrm{Ob}\, \mathcal{X}$,
\[
\Gamma\bigl( \mathcal{H}om(U , C^*) \bigr) = \Gamma(U , C^*_U)
\]
est \uncertain{une $k(U)$-algèbre}
\note{au-dessus de la ligne : « $C^*(U) =$ ».}
$k(U) = \struck{\Gamma(\ill{})}$ \ill{} \uncertain{compatible} \ill{} $=$ \ill{}
\uncertain{De même} \ill{} $\mathrm{H}^* = \mathrm{H}^*(C^*(U))$ \ill{}
$\mathrm{H}^*$\struck{\ill{}}$(U , C^*_U)$ \uncertain{est muni d’une structure} \uncertain{graduée} \ill{} d’algèbre \ill{}.
\ill{} $C^*$ \ill{} \uncertain{résolution} \ill{}
\[
(*) \qquad \mathrm{H}^*(C^*(U)) \longrightarrow \mathrm{H}^*(U , k_U)
\]
\ill{}
\marginal{\uncertain{En fait} \ill{} — E.1, E.2}
\note{suivent une douzaine de lignes illisibles, puis un diagramme :}
\[
\begin{array}{ccccccc}
\cdots \to & C^0 & \to & C^1 & \to \cdots \to & C^i & \to \cdots \\
& \| & & & & \| & \\
0 \to & \Phi & \to & \Lambda^2 \Phi & \to \cdots \to & \Lambda^i \Phi & \to \cdots
\end{array}
\]
\note{les flèches de la ligne du bas sont étiquetées $T_\wedge$.}
\ill{} \uncertain{résolution} \ill{}
\[
Z^i \simeq \Lambda^i \Psi
\]
(\ill{})
(\ill{}) \ill{}
$\cdots \longrightarrow \Lambda^{i} \longrightarrow \Lambda^{i+1}$ \ill{}

\page{351}
\note{crayon très pâle ; seules quelques formules sont lisibles.}

\struck{$Z(C^*) = \Lambda\Psi \otimes \Lambda^{i+1}$\ill{}} \quad \ill{} $C^{i+1} = \Lambda^{i+1} \Phi$

\ill{} $Z^i(C^*) = \Lambda^i \Psi$ \ill{}
\[
\begin{array}{ccc}
Z^{i+1} \otimes C^j & \xrightarrow{\;\text{donné}\;} & Z^{i+j+1} \\
\uparrow & & \uparrow \\
C^i \otimes Z^j & \xrightarrow{\;\sigma_{ij}\;} & C^{i+j} \\
\uparrow & & \\
Z^i \otimes Z^j & \xrightarrow{\;\text{donné}\;} & Z^{i+j}
\end{array}
\]
\ill{} \uncertain{trouver} $\sigma_{ij}$ \uncertain{donné} \ill{}
\[
\mathrm{Hom}_k\bigl( C^{i} \otimes Z^j , C^{i+j} \bigr) \simeq \Lambda^i \Psi \otimes \Lambda^j \Psi
\]
\ill{}
\[
\simeq \Lambda^{i+j} \Psi \oplus ( \cdots )
\]
(\uncertain{quelques termes d’homotopie} \ill{})

\section{Catégories de foncteurs et complexes de chaînes}

\page{354}
\note{la moitié supérieure du feuillet est écrite droite, la moitié inférieure tournée d’un quart de tour ; on transcrit le haut, puis le bas.}

\[
\mathcal{L} : \quad
\begin{cases}
\underline{\mathrm{Hom}}^!_{\mathcal{L}}((Ss)^\circ , \mathcal{B}) \simeq \underline{\mathrm{Hom}}^!_{\mathcal{L}}(\mathcal{B}_0 , \mathcal{B}) \\
\underline{\mathrm{Hom}}^!_{\mathcal{L}}((Ss) , \mathcal{A}) \simeq \underline{\mathrm{Hom}}^!_{\mathcal{L}}(\mathcal{A}_0 , \mathcal{A})
\end{cases}
\qquad \text{\struck{$\underline{\mathrm{Hom}}($}}
\]
\[
(Ss)^\circ \xrightarrow{\;F_0\;} \mathcal{B}_0 \simeq (\mathcal{A}b_*)^\circ \simeq (Ab^*)^\circ ,
\qquad
(Ss) \xrightarrow{\;G_0\;} \mathcal{B}
\]
\[
\mathbb{C}_* \in \quad
\begin{cases}
\mathcal{A}_0 = \mathcal{B}_0^\circ \simeq \mathcal{A}b_* , \quad \mathcal{B}_0 = \mathcal{A}_0^\circ \simeq (\mathcal{A}b^*)^\circ \\
G_0 = (F_0)^\circ , \quad F_0 = (G_0)^\circ
\end{cases}
\qquad F_0 \in \mathcal{A}b^*_*
\]
\[
\text{\struck{$\Omega_{\mathcal{L}}$}} \quad
\begin{cases}
\Omega_{\mathcal{L}} \simeq \mathcal{B}_0 \longrightarrow \mathcal{A}b_*^\circ \simeq \mathcal{A}b_\bullet^\circ \\
\Omega'_{\mathcal{L}} \simeq \mathcal{A}_0 \longrightarrow \mathcal{A}b_* \simeq \mathcal{A}b_\bullet
\end{cases}
\]
\note{les deux flèches portent une étiquette \ill{} ; la seconde est à peine tracée.}
\[
\mathbb{C}^{*\vee} \qquad
\underline{\mathrm{Hom}}(\Delta^\circ , \mathcal{A}b_*^\circ) \simeq (Ab^*_*)^\circ
\]
\note{devant $(Ab^*_*)^\circ$, un mot biffé, illisible.}
$\Omega$, $F$, $L_*$, $L_\bullet$. — \uncertain{Cas} $\mathcal{A} = \mathcal{A}b$, \ill{} \uncertain{contravariants} \ill{} \uncertain{dans} $\mathcal{A}b$.
\note{des traits relient $\Omega$, $F$, $L_*$, $L_\bullet$ à un symbole surchargé et entouré au-dessus d’eux ; un sigle est biffé à cet endroit.}
\[
\boxed{\, F(X_*) = \Omega_*(\mathbb{C}_*(X)) = \Omega_\bullet(\mathbb{C}_\bullet(X)) \,}
\]
$L_\bullet \quad L_*$
\[
\begin{aligned}
\underline{\mathrm{Hom}}_{\hat{\Delta}}(X_* , L_*) &\simeq \underline{\mathrm{Hom}}(\mathbb{C}_*(X_*) , L_*) \simeq \underline{\mathrm{Hom}}_{\mathrm{Ens}}(\mathbb{C}_\bullet(X) , L_\bullet) \\
&\simeq \underline{\mathrm{Hom}}^*(L^*_* , \mathbb{C}^*(X)) \simeq \underline{\mathrm{Hom}}^\bullet(L^\bullet , \mathbb{C}^\bullet(X))
\end{aligned}
\]
\uncertain{dans ce cas} $L_*$ \ill{}.
\note{l’indice « Ens » du troisième $\underline{\mathrm{Hom}}$ est douteux.}

Partie tournée :
\[
\mathcal{A} \subset \hat{\mathcal{A}}
\]
\[
\underline{\mathrm{Hom}}(\hat{\Delta}_{ss} \times \mathcal{A}^\circ , \mathcal{A}b) \simeq \underline{\mathrm{Hom}}(\mathcal{A}^\circ , \ldots
\]
\ill{}
\uncertain{rep.} \ill{}.
\[
\varphi \in \underline{\mathrm{Hom}}_!(\hat{\Delta}_{ss} , \hat{\mathcal{A}}) \simeq \underline{\mathrm{Hom}}(\hat{\Delta}_{ss} \times \mathcal{A}^\circ , \mathcal{A}b)
\]
\[
\mathrm{Hom}(X , \varphi(M)) = \Phi(M , X)
\]
\uncertain{rep.\ en $X$} \ill{} $M$
\struck{$= \mathrm{Hom}(M ,$} \quad $\Phi$ \quad $\underline{\mathrm{Hom}}(\mathcal{A}^\circ , \underline{\mathrm{Hom}}_!($\ill{} \quad $\hat{\Delta} \longrightarrow \mathcal{A}b$
\[
\underline{\mathrm{Hom}}_!(\hat{\Delta}_{ss} , \mathcal{A}) \longrightarrow \underline{\mathrm{Hom}}_!(\hat{\Delta} , \mathcal{A})
\ \overset{\wr}{\longrightarrow} \ \underline{\mathrm{Hom}}(\Delta , \mathcal{A})
\]
\note{le signe $\wr$ sur la dernière flèche est douteux.}
\note{une flèche courbe revient de $\underline{\mathrm{Hom}}(\Delta , \mathcal{A})$ vers $\underline{\mathrm{Hom}}_!(\hat{\Delta}_{ss} , \mathcal{A})$ ; une formule « $X = \varinjlim \ldots$ » est biffée en dessous.}
\[
\mathrm{Hom}(\varphi(X) , M) = \Phi(X , M) : \hat{\Delta}_{ss}^\circ \times \mathcal{A} \longrightarrow \mathcal{A}b
\]
(\uncertain{rep.\ en $M$} \ill{} \uncertain{dès que} \ill{} \uncertain{rep.\ en $X$})
\[
\simeq \mathrm{Hom}_{\hat{\Delta}_{ss}}(X , \psi(M)) ,
\qquad
\hat{\Delta}_{ss} \overset{\varphi}{\underset{\psi}{\rightleftarrows}} \mathcal{A}
\]
\[
\mathcal{A} \longrightarrow \ ; \qquad \mathcal{A} \times \Delta \longrightarrow \mathcal{A}b
\]

\section{Théorème d’Eilenberg--Zilber--Cartier}

\page{356}
\note{en haut à gauche « EZ 1 » ; en haut à droite un « EZ 2 » biffé. Feuillet en largeur, en deux colonnes.}

\textbf{Th.\ Eilenberg--Zilber--Cartier}

$X_{**}$ un \struck{groupe} \add{objet} abélien bisimplicial (\ill{}).
\note{une ligne biffée suit.}
\[
N \colon A_{**} \xrightarrow[\;\sim\;]{\text{éq.\ de cat}} A_{\bullet\bullet} ,
\qquad
N \colon A_* \to A_\bullet , \quad {}^\sim \colon A_* \to A_\bullet
\]
(complexe associé) $\ {}^\sim \colon A_{**} \longrightarrow A_{\bullet\bullet}$
\[
N \hookrightarrow (\,)^\sim \quad \text{q-iso} , \qquad D(\,) \ \text{dégén.\ homot.\ triv.}
\]
\note{« dégén.\ homot.\ triv. » : lecture douteuse.}
\[
\Delta \colon A_{**} \longrightarrow A_* , \qquad \int \colon A_{\bullet\bullet} \longrightarrow A_\bullet
\]
$\exists$ des \struck{quasi-iso} \add{morphismes} « naturels » fonctoriels,
\[
\int N \ \overset{\psi}{\underset{\varphi}{\rightleftarrows}} \ N \Delta
\]
qui coïncident avec les flèches évidentes en degré 0.
\note{au-dessus de « qui » un sigle, \uncertain{$X_{0,0}$}.}

\begin{tikzcd}
A_{**} \arrow[r, "N"] \arrow[d, "\Delta"'] & A_{\bullet\bullet} \arrow[d, "\int"] \\
A_* \arrow[r, "N"] & A_\bullet
\end{tikzcd}

\[
\varphi\psi \sim 1_{\int N} , \qquad \psi\varphi \sim 1_{N\Delta}
\]
\note{une flèche barrée prolonge la ligne du haut ; au carré, une petite flèche courbe marque la non-commutation stricte.}
\[
\psi\varphi \colon \int (N X_{**}) \longrightarrow \int N X_{**} ,
\qquad
\varphi\psi \colon N\Delta X_{**} \longrightarrow N\Delta X_{**}
\]
\[
1 - \psi\varphi = dk + kd , \qquad k \colon \int N X_{**} \xrightarrow{\;+1\;} \int N X_{**}
\]
\[
1 - \varphi\psi = dh + hd , \qquad h \colon N\Delta X_{**} \xrightarrow{\;+1\;} N\Delta X_{**}
\]
\note{les deux premières lignes de ce bloc écrivent $\psi\varphi$ et $\varphi\psi$ là où les homotopies placent $\varphi\psi$ et $\psi\varphi$ ; transcrit tel quel.}

\note{suit un carré où $\int N \leftrightarrows N\Delta$ et $\int {}^\sim \leftrightarrows {}^\sim\Delta$ sont reliés par des flèches verticales.}
\[
\int N \ \overset{\psi}{\underset{\varphi}{\leftrightarrows}} \ N\Delta ,
\qquad
\int {}^\sim \ \overset{\tilde{\psi}}{\underset{\tilde{\varphi}}{\rightleftarrows}} \ {}^\sim\Delta ,
\qquad
\int {}^\sim \ \overset{\tilde{\psi}}{\underset{\tilde{\varphi}}{\leftrightarrows}} \ {}^\sim(\Delta)
\]
Si $\varphi$ et $\psi$ sont choisis, pour un autre choix de $h'$, $k'$
\[
k' - k = dl + ld , \qquad h' - h = dm + md , \quad \text{etc.}\ \ldots \quad \text{$l$ et $m$ fonct.}
\]
\struck{$k - h'$}

Colonne de droite.

$\exists$ des choix traditionnels : « application d’Alexander--Whitney », « shuffles ». — cf.\ MacLane, \emph{Homology}.

\textbf{Shuffles.}
\[
\Bigl( \int \widetilde{X}_{**} \Bigr)_n = \bigoplus_{i+j=n} X_{(i,j)} \xrightarrow{\;\text{shuffles}\;} \widetilde{X}_{(n,n)}
\]
\[
x_{ij} \longmapsto x_{i+j,\,i+j}
\]
\note{entre les deux membres, une somme biffée, illisible.}
$(i,j)$-shuffle : $(1 , \ldots , i+j = n) \mapsto s(1) \ldots s(i) \ \ s(i+j)$
\[
x_{i1} \qquad x_{i+1,\,i+1}
\]
\[
\sum_{\substack{0 \le \alpha \le i \\ \mathrm{sgn}(\sigma)}} (\pm)\, s'_0 s'_1 \cdots \widehat{s_\alpha} \cdots s'_i \, x_{i1}
\qquad
\sum_{\text{sur tout } \sigma \in \Sigma_n} s_{\sigma(i+1)} \cdots s_{\sigma(i+j)} \, s'_{\sigma(1)} \cdots s'_{\sigma(i)}
\]
\note{le coefficient noté $(\pm)$ est illisible ; seul « sgn$(\sigma)$ » se lit dessous.}
t.q.\ $\sigma(\alpha) < \sigma(\beta)$ si $\alpha < \beta \le i$ ou \struck{\ill{}} $i < \alpha < \beta \le n$.
\drawing{sous deux crochets illisibles, un faisceau de segments reliant deux colonnes de points : l’image d’un shuffle.}

$A \qquad \overline{B}(A)$

\textbf{Alexander--Whitney}
\[
\widehat{X}_{n,n} \longrightarrow \bigoplus X_{ij} ,
\qquad
x_{n,n} \longmapsto \sum_{i=0}^{n} (-1)^{i} \, d_0^{\,n-i} \, \tilde{d}^{\,i} x_{n,n}
\]
\note{le signe $(-1)^i$ est de lecture douteuse.}
$\tilde{d}$ signifie le dernier op.\ face, $d'$ \uncertain{seconde} op.\ simpl.
\[
\tilde{d}^{\,i} = d^{\prime n} d^{\prime n-1} \cdots d^{\prime n-i+1}
\]
\[
C^i(X , A) \otimes C^j(X , B) \longrightarrow C^{i+j}(X \times X , A \otimes B) \longrightarrow C^{i+j}(X , A \otimes B)
\]
\[
f , \ g ; \qquad \mathbb{Z}X \otimes \mathbb{Z}X \longleftarrow \mathbb{Z}(X \times X)
\]
E--Z--C $\Rightarrow$ E--Z.

Cup-produit :
\[
h(x_1 , \ldots , x_{n+1}) = \sum (-1)^{\cdots} f(x_1 , \ldots , x_i)\, g(x_{i+1} , \ldots , x_n)
\]
\note{l’exposant du signe est illisible.}
associatif (non commutatif, mais comm.\ à homotopie près) $\Rightarrow$ $\mathrm{H}^*(X , R)$ est anticomm.

\page{357}
\[
K(\mathbb{Z} , i) \times K(\mathbb{Z} , j) \longrightarrow K(\mathbb{Z} , i+j)
\]
\[
\searrow \quad K(\mathbb{Z} , i) \otimes K(\mathbb{Z} , j) \xrightarrow{\;f\;} K(\mathbb{Z} , i+j)
\]
Il suffit de construire $Nf$
\[
N\bigl( K(\mathbb{Z} , i) \otimes K(\mathbb{Z} , j) \bigr)
\xrightarrow{\;\mathrm{E\text{-}Z}\;}
N K(\mathbb{Z} , i) \otimes N K(\mathbb{Z} , j)
\longrightarrow N K(\mathbb{Z} , i+j) = \mathbb{Z}[i+j]
\]
\[
N K(\mathbb{Z} , i) \otimes N K(\mathbb{Z} , j) = \mathbb{Z}[i] \otimes \mathbb{Z}[j]
\]
\struck{$\bigl( \mathrm{Inv}^G(S\mathfrak{g}^*) , \ \mathrm{H}^*(G , \mathrm{L}\,\mathrm{Sym}\, \mathfrak{g}^*) \bigr) \Rightarrow \mathrm{H}^*(BG)$}
\note{ces trois termes, mis entre parenthèses, sont barrés d’une grande croix.}

\page{358}
\note{en haut à droite « EZ 2 ». Feuillet en largeur, en deux colonnes.}

Cor. $\Sigma_n \to E\Sigma_n =$ \struck{\ill{}} \quad \uncertain{$\Sigma_n \times \Sigma_n = \Sigma_n$}
\[
\downarrow \qquad K(\Sigma_n , 1)
\]
\[
W_n = \mathbb{Z} E\Sigma_n \xrightarrow{\;\sim\;} \mathbb{Z}
\]
résol.\ de $\mathbb{Z}$ par des $\Sigma_n$-modules \uncertain{plats} libres.
Cor.\ du th.
\[
W_n \xrightarrow{\;\sim\;} \mathrm{Hom}^\bullet\bigl( \mathbb{Z}(X \times X \times \cdots) , \mathbb{Z}X \otimes \cdots \otimes \mathbb{Z}X \bigr)
\]
\note{devant $\mathbb{Z}(X \times X \times \cdots)$, un premier essai est biffé.}
D’où par adjonction une flèche
\[
W_n \otimes \mathbb{Z}(X \times \cdots \times X) \xrightarrow{\;\lambda\;} \mathbb{Z}X \otimes \mathbb{Z}X \otimes \cdots \otimes \mathbb{Z}X
\]
\[
(\text{resp.\ } W_n \otimes \mathbb{Z}X \otimes \cdots \otimes \mathbb{Z}X \longrightarrow \mathbb{Z}(X \times \cdots \times X))
\]
Ceci exprime l’homotopie-commutativité \uncertain{annoncée} \uncertain{plus haut}.

\emph{Ex.\ $n = 2$}, 1) exemple topologique :
\[
X \times X \xrightarrow{\;m\;} X \quad \text{homotopiquement comm.}
\]
\[
\Sigma_2\text{-équiv.} \quad I \times X \times X \longrightarrow X \quad \cdots
\]
\note{la ligne $I \times X \times X \to X$ est entourée, et une flèche pointillée en part vers la droite.}
plus fort :
\[
E\Sigma_2 \times X^2 \longrightarrow X \qquad \Sigma_2\text{-équiv.}
\]
Classiquement
\[
E\Sigma_n \times X^n \longrightarrow X \qquad \text{relevant la mult.\ itérée}
\]
$\Sigma_2$ agit sur $S^n$ par action antipodale, $S^n \to \mathbb{P}^n(\mathbb{R})$ ;
\[
E\Sigma_2 = \varinjlim S^n
\]
\[
S^1 \times X^2 \longrightarrow X \qquad x , y \qquad \text{\struck{$xy = yx$}} \quad x y x^{-1} y^{-1}
\]
\[
S^2 \times X^2 \xrightarrow{\;\Sigma_2\text{-équiv.}\;} X , \qquad \text{etc.}\ \ldots
\]
\drawing{une petite sphère hachurée, cerclée d’un équateur épais.}

Colonne de droite.
\[
E(\Sigma_n \wr \Sigma_j) \times (X_1^0 \cdots X_j^0) \times X_1^1 \cdots X_j^1 \times \cdots \times X_1^n \times \cdots \times X_j^n
\]
\[
\downarrow \qquad ?
\]
\[
X \longleftarrow E\Sigma_n \times X^n
\]
\note{à gauche de la première ligne, une flèche vient d’un $E\Sigma_n \times (E\Sigma_j \times X)^n$ biffé, d’où descend une flèche vers $X^n$.}

Cor. : les énoncés précédents impliquent que lorsque $X$ est un H-esp.\ satisfaisant à ce genre de cdn de commutativité à homotopie près, alors les chaînes $\mathbb{Z}X$ sont \uncertain{des}
\[
W_n \otimes \mathbb{Z}X^{\otimes n} \xrightarrow{\;\Sigma_n\text{-équiv.}\;} \mathbb{Z}X^{\otimes} \longrightarrow
\]
\note{l’exposant de $\mathbb{Z}X$ au but est mal formé.}
t.q.\ induise la mult.\ des chaînes
\[
\mathbb{Z}X^{\otimes n} \xrightarrow[\;\text{shuffle}\;]{\;\mathrm{EZ}\;} \mathbb{Z}(X^n) \xrightarrow{\;\mathbb{Z}(m)\;} \mathbb{Z}X
\]

Constr.\ des opérations de Steenrod : $A$ un anneau de car.\ $p$. \struck{$\mathbb{Z}X \longrightarrow$}
\[
f \in \mathrm{H}^n(X , A) , \qquad \widetilde{\mathbb{Z}X} \xrightarrow{\;f\;} A[n]
\]
\[
W_p \otimes \mathbb{Z}(X^p) \xrightarrow{\;\lambda\;} \widetilde{\mathbb{Z}X}^{\otimes p} \xrightarrow{\;f^{\otimes p}\;} A[n]^{\otimes p} \xrightarrow{\;\mu\;} A[pn]
\]
\[
\nearrow \Delta \qquad W_p \otimes \mathbb{Z}X \quad \Sigma_n\text{-équiv.}\ (\text{avec action triv.\ sur } \mathbb{Z}X)
\]
\marginal{encadré : cf.\ Steenrod, \emph{Coh.\ ops}, p.~99}
\[
Pf = \mu \circ f^{\otimes p} \circ \lambda \in {}_{\Sigma_n}\mathrm{H}^{pn}(\mathbb{Z}(X^p) , A)
\]
\[
\downarrow
\]
\[
(W_n \otimes \mathbb{Z}X)/\Sigma_n = \mathbb{Z}X \otimes \mathbb{Z}K(\Sigma_n , 1) ,
\qquad
{}_{\Sigma_n}\mathrm{H}^{pn}(X , A)
\]
\[
\| \qquad \mathrm{H}^*(K(\Sigma_p) \times X , A)
\]
\[
\text{Künneth} \quad \| \qquad \mathrm{H}^*(\Sigma_p) \otimes_{\mathbb{F}_p} \mathrm{H}^*(X , A)
\]

\page{359}
\note{en haut à droite « EZ 3 ».}

Les applications shuffles et A-W sont \emph{associatives}
\begin{itemize}
\item et \emph{shuffles} est (anti-)commutative
\item A-W est seulement comm.\ à homotopie près.
\end{itemize}
On peut exprimer ceci avec des diagrammes d’objets trisimpliciaux.
\marginal{Réf.\ Illusie, \emph{Thèse}.}

\begin{tikzcd}
X_{***} \arrow[r, "N_{**}"] \arrow[d, "\Delta_{12}"'] & X_{\bullet\bullet\bullet} \arrow[r, "\int_{12}"] & X_{\bullet\bullet} \arrow[d, "\int"] \\
X_{**} \arrow[r, "N"] & X_{\bullet\bullet} & X_\bullet
\end{tikzcd}

Clair pour $\mathbb{Z}X$.
\struck{\ill{}}

L’énoncé de commutativité s’énonce comme suit. $X$ ens.\ simpl.
\struck{$\mathrm{Hom}(\mathbb{Z}X \otimes \mathbb{Z}X , \mathbb{Z}(X \times X))$}
\[
\mathrm{Hom}^\bullet\bigl( \mathbb{Z}(X \times X)^\sim , \mathbb{Z}\widetilde{X} \otimes \mathbb{Z}\widetilde{X} \bigr) \longrightarrow \mathbb{Z}
\]
\[
\mathrm{Hom}\bigl( \mathbb{Z}(X_0 \times X_0) \xrightarrow{\;\sim\;} \mathbb{Z}X_0 \otimes \mathbb{Z}X_0 \bigr) \longrightarrow \mathbb{Z}
\]
\marginal{$\Sigma_2$ agit sur ce complexe}
\[
X \longrightarrow \mathrm{pt}
\]
\[
\mathrm{Hom}\bigl( \mathbb{Z}(\ ,\ ) , \widetilde{\mathbb{Z}}(\ ) \otimes \widetilde{\mathbb{Z}}(\ ) \bigr) \longrightarrow \mathbb{Z}
\]

\textbf{Théorème} (Dold) : « Über Steenrodsche Operationen … »

Pour tout $X$, ce complexe
\[
\mathrm{Hom}^\bullet\bigl( \mathbb{Z}(X \times \cdots \times X) , \mathbb{Z}(X) \otimes \cdots \otimes \mathbb{Z}X \bigr)
\]
\[
\mathrm{Hom}^\bullet\bigl( \mathbb{Z}X \otimes \cdots \otimes \mathbb{Z}X , \mathbb{Z}(X \times \cdots \times X) \bigr)
\]
est une résolution libre de $\mathbb{Z}$ dans la cat.\ des $\Sigma_n$-modules pour $X$ variable.

On bricole pour avoir un complexe de chaînes

\page{360}
\note{en haut à droite « EZ 3 », le chiffre repassé, peut-être sur un autre.}
\[
f \in \mathrm{H}^n(X , A) \Longrightarrow Pf \in {}_{\Sigma_n}\mathrm{H}^{pn}(E\Sigma_n \times X^p , A)
\Longrightarrow \Delta Pf \in \mathrm{H}^{pn}(K(\Sigma_p , 1) \times X , A)
\]
\[
\mathrm{H}^{pn}(K(\Sigma_p , 1) \times X , A) = \mathrm{H}^*(K(\Sigma_p , 1)) \otimes_{\mathbb{F}_p} \mathrm{H}^*(X , A)
\]
\struck{\ill{}}
\[
e_\alpha \in \mathrm{H}_\alpha(K(\Sigma_p , 1))
\qquad \searrow \qquad
D_i x \in \mathrm{H}^{pn-i}(X , A)
\]
\[
P^i x = (-1)^{\cdots} \, D_{f(i)}
\]
\note{l’exposant du signe est illisible.}
Ce reindexage est ce qui donne la formule de Cartan.

En fait on considère $\mathbb{Z}_p \subset \Sigma_p$ et on travaille avec la coh.\ du groupe cyclique.
\note{l’argument s’interrompt ici ; il se poursuit, s’il se poursuit, au lot suivant.}

\end{document}
